\section{Preliminaries}\label{sec:prelim}

Throughout, $\A=\{0,1\}^2$, and a cell $(x,y)\in\A$ is also written as the integer $2x+y\in\{0,1,2,3\}$; thus $3=(1,1)$ and $0=(0,0)$. The left shift is $\sh(c)_i=c_{i+1}$. A \emph{block map of radius $R$} is a map $g\colon\A^{\Z}\to\A^{\Z}$ of the form $g(c)_i=\gamma(c_{i-R},\dots,c_{i+R})$ for a local rule $\gamma\colon\A^{2R+1}\to\A$. A block map of radius $R$ is also a block map of every larger radius, so ``radius at most $R$'' and ``radius $R$'' describe the same set of maps. By the Curtis--Hedlund--Lyndon theorem \cite{Hedlund1969}, the continuous maps of $\A^{\Z}$ into itself that commute with $\sh$ are the block maps, and $\Aut(\A^{\Z})$ is the group of invertible ones.

For an elementary rule $f$, identified with its Wolfram number $\sum_{a,b,c}f(a,b,c)2^{4a+2b+c}$, the automaton $F_f$ is defined by \eqref{eq:F}. It has radius $1$, and $F_f^b$ has radius at most $|b|$, so $\sh^aF_f^b$ is a block map of radius at most $|a|+|b|$ that commutes with $F_f$.

\begin{definition}
The \emph{spacetime subshift} of a reversible cellular automaton $F$ on $\A^{\Z}$ is
\[
X_F=\bigl\{(F^n(c))_{n\in\Z}:\ c\in\A^{\Z}\bigr\}\subset\A^{\Z^2},
\]
with the $\Z^2$-action generated by the horizontal shift $\sh$ and the vertical shift $\tau$, $(\tau z)_{n}=z_{n+1}$. An \emph{automorphism} of $X_F$ is a homeomorphism of $X_F$ commuting with both.
\end{definition}

\begin{lemma}\label{lem:spacetime}
For reversible $F$ the map $c\mapsto(F^n(c))_{n\in\Z}$ is a homeomorphism $\A^{\Z}\to X_F$ conjugating $(\sh,F)$ to $(\sh,\tau)$, and it induces an isomorphism
\[
\Aut(X_F)\;\cong\;\{g\in\Aut(\A^{\Z}):\ gF=Fg\}.
\]
Under this isomorphism $\langle\sh,\tau\rangle$ corresponds to $\langle\sh,F\rangle$.
\end{lemma}

\begin{proof}
The map is continuous, injective (the $0$-th row recovers $c$), and surjective onto $X_F$ by definition, hence a homeomorphism of compact spaces; it carries $\sh$ to $\sh$ and $F$ to $\tau$ since $F^{n}(F(c))=F^{n+1}(c)$. A homeomorphism of $\A^{\Z}$ commuting with $\sh$ is a block map by Curtis--Hedlund--Lyndon, and commuting with $\tau$ means commuting with $F$.
\end{proof}

Thus determining $\Aut(X_F)$ means determining which invertible block maps commute with $F$, and we regard the elements of $\Aut(X_F)$ as block maps. The block maps that commute with $F$ and are not invertible correspond in the same way to the endomorphisms of $X_F$ that are not automorphisms; if $F$ fixes a constant configuration, the constant map onto it is one. A constant configuration fixed by $F_f$ is of the form $(a,a)^{\Z}$ with $f(a,a,a)=0$; for $f\in\{91,95\}$ the only one is $(1,1)^{\Z}$, and for $f=94$ there are two, $(0,0)^{\Z}$ and $(1,1)^{\Z}$.

\subsection{Symmetries}
Three involutions of $\A^{\Z}$ act on second-order automata: the cellwise complement $\kappa(x,y)=(1-x,1-y)$, the reflection $\mathrm{rev}(c)_i=c_{-i}$, and the time reversal $\rho(x,y)=(y,x)$, applied in every cell. The first two are the Boolean conjugation and the reflection of \cite[Section~2]{Takesue1989}, and that exchanging the two rows reverses time is noted in \cite[p.~2500]{Takesue1987}.

\begin{lemma}\label{lem:classes}
For every elementary rule $f$, $\kappa F_f\kappa=F_{\kappa f}$ with $\kappa f(a,b,c)=f(1-a,1-b,1-c)$; $\mathrm{rev}\,F_f\,\mathrm{rev}=F_{f^r}$ with $f^r(a,b,c)=f(c,b,a)$, and $\mathrm{rev}\,\sh\,\mathrm{rev}=\sh^{-1}$; and $\rho F_f\rho=F_f^{-1}$. For $\theta\in\{\kappa,\mathrm{rev}\}$ the map $g\mapsto\theta g\theta^{-1}$ preserves the radius of block maps, maps the block maps that commute with $F_f$ onto those that commute with $F_{f'}$ and $\Aut(X_{F_f})$ onto $\Aut(X_{F_{f'}})$, where $f'=\kappa f$ or $f^r$, and maps $\langle\sh,F_f\rangle$ onto $\langle\sh,F_{f'}\rangle$.
\end{lemma}

\begin{proof}
By \eqref{eq:F}, $F_f(\kappa c)_i=(1-y_i,\,1-x_i\oplus f(1-y_{i-1},1-y_i,1-y_{i+1}))$, and applying $\kappa$ gives $(y_i,\,x_i\oplus\kappa f(y_{i-1},y_i,y_{i+1}))$. Similarly $\mathrm{rev}\,F_f\,\mathrm{rev}(c)_i=F_f(\mathrm{rev}\,c)_{-i}=(y_i,\,x_i\oplus f(y_{i+1},y_i,y_{i-1}))$. And $F_f(\rho c)_i=(x_i,\,y_i\oplus f(x_{i-1},x_i,x_{i+1}))$, whose image under $\rho$ is $F_f^{-1}(c)_i$. Conjugation by a cellwise map or by $\mathrm{rev}$ preserves the radius, and the remaining statements follow, since $\kappa$ commutes with $\sh$ and $\mathrm{rev}$ conjugates $\sh$ to $\sh^{-1}$.
\end{proof}

So $\Aut(X_{F_f})$, and every statement below about block maps of bounded radius, depends only on the class of $f$ under $\kappa$ and reflection; there are $88$ classes \cite{Takesue1987}, \cite[Table~I]{Takesue1989}. The rules $91$, $94$ and $95$ are invariant under reflection, and their classes are $\{91,218\}$, $\{94,122\}$ and $\{95,250\}$; the class of $155$ is $\{155,203,211,217\}$.

\subsection{Diagrams}
For $c\in\A^{\Z}$ the \emph{diagram} of $c$ is $Y_c\colon\Z^2\to\{0,1\}$, where $Y_c(t,i)$ is the $y$-bit of $F^t(c)_i$. The $x$-row of $F(d)$ is the $y$-row of $d$, so the $x$-bit of $F^t(c)_i$ is $Y_c(t-1,i)$, and $c_i=(Y_c(-1,i),Y_c(0,i))$. By \eqref{eq:F} every diagram satisfies the \emph{spacetime relation}
\begin{equation}\label{eq:rel}
Y(t+1,i)\oplus Y(t-1,i)=f\bigl(Y(t,i-1),Y(t,i),Y(t,i+1)\bigr)\qquad\text{for all }(t,i)\in\Z^2,
\end{equation}
so two consecutive rows determine all rows, and $c\mapsto Y_c$ is a bijection from $\A^{\Z}$ onto the solutions of \eqref{eq:rel}. For $t\ge0$, $Y_c(t,i)$ depends only on $c_{[i-t,i+t]}$. Moreover $Y_{\sh c}(t,i)=Y_c(t,i+1)$, $Y_{F(c)}(t,i)=Y_c(t+1,i)$ and $Y_{\rho c}(t,i)=Y_c(-1-t,i)$, the last because $F^t(\rho c)=\rho F^{-t}(c)$.

\subsection{The fixed points of \texorpdfstring{$\sh^{\pm1}F_f$}{sigma F}}\label{sec:fixed}
The fixed points of $\sh F_f$ and of $\sh^{-1}F_f$ are the configurations whose diagrams are constant along the light rays of one of the two families. They are described by a shift of finite type on $\{0,1\}$.

\begin{lemma}\label{cl:lem:Z2}
For an elementary rule $f$ let $\Phi_f\subseteq\{0,1\}^{\Z}$ be the set of sequences $\varphi$ with
\[
\varphi(k-1)\oplus\varphi(k+1)=f\bigl(\varphi(k-1),\varphi(k),\varphi(k+1)\bigr)\qquad\text{for all }k\in\Z.
\]
The maps $\varphi\mapsto c$, $c_i=(\varphi(i+1),\varphi(i))$, and $\varphi\mapsto d$, $d_i=(\varphi(i-1),\varphi(i))$, are bijections from $\Phi_f$ onto the fixed points of $\sh F_f$ and of $\sh^{-1}F_f$ respectively. If $\Phi_f$ contains a sequence that is not periodic, then $(a,b)\mapsto\sh^aF_f^b$ is injective, so $\langle\sh,F_f\rangle\cong\Z^2$ and $F_f$ has infinite order.
\end{lemma}

\begin{proof}
If $c_i=(\varphi(i+1),\varphi(i))$, then by \eqref{eq:F}, $F_f(c)_i=\bigl(\varphi(i),\,\varphi(i+1)\oplus f(\varphi(i-1),\varphi(i),\varphi(i+1))\bigr)$, which is $(\varphi(i),\varphi(i-1))=c_{i-1}$ exactly when the displayed equation holds at $k=i$. So $F_fc=\sh^{-1}c$, that is $\sh F_fc=c$. Conversely, if $\sh F_fc=c$, then $\varphi(i)=y(c_i)$ satisfies $x(c_i)=y(F_f^{-1}c)_i=y(\sh c)_i=\varphi(i+1)$, and the equation holds by the same computation. The case of $\sh^{-1}F_f$ is identical with $F_fd=\sh d$. Let $\varphi\in\Phi_f$ be non-periodic, so that the corresponding $c$ and $d$ are not periodic under $\sh$. If $\sh^aF_f^b=\mathrm{id}$, then $c=\sh^aF_f^bc=\sh^{a-b}c$ and $d=\sh^{a+b}d$, so $a-b=a+b=0$.
\end{proof}

The set $\Phi_f$ is the shift of finite type on $\{0,1\}$ given by the allowed triples $(u,v,w)$ with $u\oplus w=f(u,v,w)$. For the four rules of this paper it is the set of sequences that contain no $000$ for $f=91$; no $010$ for $f=94$; neither $000$ nor $010$, that is, no two zeros at distance $2$, for $f=95$; and none of $000$, $110$, $111$ for $f=155$ (Sections~\ref{cl:sec:light} and~\ref{cl:sec:155}). In each case $\Phi_f$ has positive entropy and contains sequences that are not periodic, for instance $\cdots111\,0\,111\cdots$ for $91$ and $95$, $\cdots000\,111\cdots$ for $94$ and $\cdots0101\,001\,0101\cdots$ for $155$. So by Lemma~\ref{cl:lem:Z2},
\[
\langle\sh,F_f\rangle\cong\Z^2\qquad\text{for }f\in\{91,94,95,155\}.
\]
For rule $110$, by contrast, $\Phi_{110}$ consists of $0^{\Z}$, $1^{\Z}$ and one $\sh$-orbit of non-periodic points, the \emph{kinks} on which the determination of $\Aut(X_{F_{110}})$ in \cite{PaperA} is built.

\subsection{Blocking words, window graphs and eigenfields}\label{sec:blocking}
Let $F$ be a cellular automaton of radius $1$ on $\A^{\Z}$. A word $w$ on an interval $I$ of cells is \emph{blocking} for $F$ (cf.\ the proof of \cite[Theorem~3]{Kurka1997} and, for the term, \cite[Section~2.2]{BT2000}) if there is a cell $i_0\in I$ such that $F^n(c)_{i_0}=F^n(c')_{i_0}$ for all $n\ge0$ and all configurations $c,c'$ that carry $w$ on $I$. A word that is fixed by $F$ in every context is blocking.

Let $F$ be a bijective cellular automaton on $\A^{\Z}$. For $(k,p)\in\Z^2$ with $p\ge1$ and $M\ge0$, let $W^{(k,p)}_M$ be the graph on $\A^{M+1}$ with an edge from each word to each of its possible images under $\sh^kF^p$, the cells outside the window ranging over all values: there is an edge from $c_{[0,M]}$ to $\sh^kF^p(c)_{[0,M]}$ for every $c\in\A^{\Z}$. The graph $W^{(1,1)}_M$ is the light-speed window graph used for the permutive rules in \cite{PaperB}, and $W^{(0,1)}_M$ is a zero-speed graph. A nonconstant block map $\vartheta\colon\A^{\Z}\to B^{\Z}$ to sequences over a finite alphabet $B$ with $\vartheta\circ F^p=\sh^{-k}\circ\vartheta$ is an \emph{eigenfield for $(k,p)$}, of speed $k/p$; composing with the indicator of a letter shows that one exists if and only if one exists with $B=\{0,1\}$.

\begin{lemma}\label{cl:lem:eigen}
Let $F$ be a bijective cellular automaton on $\A^{\Z}$ and $(k,p)\in\Z^2$ with $p\ge1$. The following are equivalent.
\begin{enumerate}
\item[(1)] The graph $W^{(k,p)}_M$ is strongly connected for every $M\ge0$.
\item[(2)] There is no nonconstant block map $\vartheta\colon\A^{\Z}\to\{0,1\}^{\Z}$ with $\vartheta\circ F^p=\sh^{-k}\circ\vartheta$.
\item[(3)] Every continuous function $v\colon\A^{\Z}\to\{0,1\}$ with $v\circ\sh^kF^p=v$ is constant.
\end{enumerate}
\end{lemma}

\begin{proof}
Write $\varsigma=\sh^kF^p$. If $\vartheta$ is as in (2), then $v(c)=\vartheta(c)_0$ is continuous and $v\circ\varsigma=v$, because $\vartheta$ commutes with $\sh$; conversely, if $v$ is as in (3), then $\vartheta(c)_i=v(\sh^ic)$ is a block map with $\vartheta\circ F^p=\sh^{-k}\circ\vartheta$, and $\vartheta$ is constant if and only if $v$ is. So (2) and (3) are equivalent.

Suppose that (3) fails, and let $v$ be a nonconstant continuous function with $v\circ\varsigma=v$. It depends on the cells of some interval, which after replacing $v$ by $v\circ\sh^a$ is $[0,M]$; let $T$ be the set of words $w\in\A^{M+1}$ with $v=1$ on the configurations that carry $w$ on $[0,M]$. Every edge of $W^{(k,p)}_M$ leads from $c_{[0,M]}$ to $\varsigma(c)_{[0,M]}$ for some $c$, and $v(\varsigma(c))=v(c)$; so no edge joins $T$ and its complement, both of which are nonempty, and (1) fails.

Suppose that (1) fails for some $M$. Every vertex of $W^{(k,p)}_M$ has an outgoing edge, so the graph has a strongly connected component $T$ that no edge leaves, and $T$ is not the whole vertex set. Let $U$ be the set of configurations $c$ with $c_{[0,M]}\in T$; it is open and closed, nonempty and not all of $\A^{\Z}$, and $\varsigma(U)\subseteq U$. The uniform Bernoulli measure on $\A^{\Z}$ is invariant under $\sh$ and under the surjective cellular automaton $F$ \cite{Hedlund1969}, so $\varsigma(U)$ and $U$ have the same measure; since $\varsigma$ is a homeomorphism, $U\setminus\varsigma(U)$ is an open set of measure $0$, hence empty. So $\varsigma(U)=U$, the indicator function of $U$ is a nonconstant continuous function invariant under $\varsigma$, and (3) fails.
\end{proof}

Condition (1) says that $\sh^kF^p$ is chain transitive. By the lemma, the graphs $W^{(\pm1,1)}_M$ are strongly connected for every $M$ if and only if $F$ has no eigenfield for $(k,p)=(\pm1,1)$, and the graphs $W^{(0,1)}_M$ if and only if it has none for $(0,1)$. An eigenfield for $(0,1)$ is a local conservation law in the sense of Takesue \cite{Takesue1987}, \cite[Section~4]{Takesue1989}, and $W^{(0,1)}_M$ is the graph of the transition matrix on words of $M+1$ cells considered there; for $(k,p)=(0,1)$ the equivalence of (1) and (2) is stated in \cite[Section~4]{Takesue1989} for the second-order elementary automata (see also \cite[Proposition~4.1]{HT1991}). The proof above is included because the lemma is needed for all $(k,p)$ and no proof is written out there.
